\chapter{Conclusão}

As divergências entre as percepções econômicas do público geral e o consenso técnico estão associadas sobretudo à posição ideológica autodeclarada, mas não exclusivamente a ela. A formação econômica avançada está associada a respostas sistematicamente distintas quando o item tem componente factual dominante; essa diferenciação desaparece quando a crença funciona como marcador de identidade. Esse é o resultado central deste trabalho — e ele deve aparecer aqui, no começo, porque a conclusão é uma resposta, não uma narrativa que amadurece lentamente até revelar sua moral.

Com 184 respondentes, 53 variáveis dependentes e dois preditores centrais, os resultados revelam duas histórias sobrepostas. A primeira pertence à ideologia: 34 das 53 DVs apresentam efeito estatisticamente significativo do espectro político após correção de Benjamini-Hochberg. Em blocos como H3 (Antimercado) e H7 (Ideologia), essa taxa chega a 9/9 e 11/12, respectivamente. A posição política autodeclarada é o preditor mais consistente, mais robusto e mais ubíquo das percepções econômicas desta amostra.

A segunda história pertence à formação econômica: cinco efeitos nominalmente significativos, nenhum sobrevivendo a BH, todos na direção correta segundo a literatura. Com poder de apenas $28{,}2\%$ para detectar OR $= 2{,}0$ com $n_{\texttt{econ}} = 17$, um resultado nulo após correção de multiplicidade era o desfecho previsível. O estudo não tem força confirmatória; tem força exploratória — e declará-lo assim não é fraqueza metodológica, mas honestidade que protege o leitor e o próprio pesquisador.

\section{O que a Dominância Ideológica Ensina}
\label{sec:dominancia-ideologica-licao}

Que a ideologia condicione crenças econômicas não é novidade. O que este estudo acrescenta é a dimensão quantitativa desse fenômeno no contexto brasileiro e a constatação de que ele atravessa transversalmente todos os temas — de juros a importações, de reformas estruturais a pessimismo sobre renda.

\citeonline{Sapienza2013} mostraram que mesmo especialistas em economia divergem de outros economistas quando suas orientações políticas diferem; a diferença entre economistas de esquerda e de direita em certas perguntas é tão grande quanto a diferença entre economistas e leigos. Os dados aqui obtidos são consistentes com esse achado: a formação econômica não blinda contra a ideologia — ela convive com ela, cada qual operando em registro próprio.

\citeonline{kahan2012polarization} chamou esse processo de \textit{cognição cultural}: indivíduos com maior capacidade cognitiva frequentemente usam essas capacidades para defender com maior sofisticação as posições que já sustentavam antes. \citeonline{taber2006motivated} documentaram o mecanismo experimental: o ceticismo é seletivo — aplicado com mais rigor a informações que ameaçam a crença prévia do que a informações que a confirmam. \citeonline{zaller1992nature} generalizou: a opinião pública é recepção e aceitação seletiva de mensagens, mediada por predisposições que funcionam como filtros permanentes.

Esse quadro tem uma consequência prática desconfortável: comunicação de política pública baseada exclusivamente em evidência tende a subestimar o problema. A ``verdade técnica'' isolada não apenas é insuficiente em certos contextos — pode ser contraproducente, ao mobilizar a defesa identitária do interlocutor. Não por má-fé, mas porque, como Dostoiévski intuiu no homem do subsolo, nem toda resistência à razão nasce da ignorância: às vezes nasce do desejo de não ser vencido por ela, de manter intacta a coerência interna de um mundo em que a crença não precisa de validação externa. Reconhecer isso não é niilismo epistemológico — é o primeiro passo para uma comunicação mais eficaz de resultados técnicos.

\section{Formação Econômica como Preditor Diferencial}
\label{sec:formacao-calibragem}

Nos itens com maior evidência de associação com a formação econômica, a componente factual domina a componente valorativa. Os três itens com maior diferença condicional são:

\begin{enumerate}[label=\roman*)]
    \item \textbf{Livre-comércio} (\textit{produtos importados benéficos}, $\Delta\bar{Y} = +0{,}37$): respondentes com pós-graduação em Economia apresentam maior endosso à abertura comercial. Esse padrão é consistente com o consenso documentado por \citeonline{Fuller2014}. A diferença condicional não pode ser atribuída causalmente à formação econômica, pois quem escolhe esse percurso acadêmico pode diferir sistematicamente do grupo de controle em dimensões não observadas.

    \item \textbf{Produtividade} (\textit{produtividade aumentando devagar}, $\Delta\bar{Y} = +0{,}60$): respondentes com pós-graduação em Economia apresentam maior preocupação com a estagnação da produtividade do que o grupo de controle. Essa inversão de sinal sugere que a formação está associada à identificação de uma fricção estrutural que o cotidiano não torna saliente — análogo ao ``o que não se vê'' de Bastiat \cite{bastiat2010oqueseve} no sentido técnico.

    \item \textbf{Déficit fiscal} (\textit{déficit federal grande}, $\Delta\bar{Y} = -0{,}46$): o grupo de controle apresenta maior concordância com a gravidade do déficit em comparação ao grupo de economistas. A diferença condicional é consistente com a literatura que distingue sustentabilidade fiscal de curto prazo de pressões intertemporais gerenciáveis \cite{franco2017moeda}.
\end{enumerate}

Em contrapartida, nos itens em que valor e identidade dominam — questões distributivas (H3), avaliação de governos (H7), reformas trabalhistas e previdenciárias — a formação econômica não está associada a diferenças discerníveis além do nível de significância alcançável com $n_{\texttt{econ}} = 17$. A implicação não é que o conhecimento seja irrelevante nesses campos, mas que não produz diferenciação associativa detectável nesta amostra quando o item mobiliza pertencimento político. O espectador imparcial que disciplina paixões por meio de normas de justiça internalizadas \cite{smith1759-theory-of-moral-sentiments} é mais fácil de observar em temas com componente factual dominante do que em questões onde a própria definição de ``justo'' é o ponto em disputa.

\section{Implicações para Pesquisa Futura: o Que os Dados Sugerem e o Que Não Provam}
\label{sec:arco-institucional}

\subsection*{O que os dados suportam}

Os padrões associativos documentados — especificamente, que posição ideológica e formação econômica avançada estão associadas a respostas sistematicamente distintas em 34 e 5 itens, respectivamente — são consistentes com a hipótese de que predisposições ideológicas e treinamento técnico moldam percepções econômicas de modo diferenciado. Essa consistência com \citeonline{Sapienza2013} e \citeonline{The_Myth_of_the_Rational_Voter} constitui evidência exploratória e sugere que o fenômeno é transferível ao contexto brasileiro.

Os dados suportam também a caracterização do problema: crenças economicamente frágeis coexistem com forte ancoragem ideológica em itens como reformas estruturais, avaliação fiscal e abertura comercial. Como \citeonline{kahan2012polarization} documentou, maior capacidade cognitiva frequentemente amplifica a defesa de posições preexistentes em vez de corrigi-las — o que sugere que intervenções puramente informativas têm alcance limitado em itens identitariamente carregados.

\subsection*{O que os dados não suportam}

Os dados deste estudo \textit{não} permitem:

\begin{enumerate}[label=\roman*)]
    \item Atribuir causalidade à formação econômica: sem variação exógena em \texttt{econ}, o estimando $\hat{\beta}_{\texttt{econ}}$ captura a diferença condicional entre grupos que se auto-selecionam, não o efeito de uma intervenção. Quem faz mestrado ou doutorado em Economia difere sistematicamente do grupo de controle em dimensões não observadas (capacidade analítica prévia, motivação intrínseca, contexto institucional) que nenhum conjunto de controles observados pode eliminar \cite{angrist2009}.
    \item Generalizar para a população brasileira: a amostragem não probabilística limita os resultados ao quadro amostral coletado. Prevalências observadas não são estimativas representativas.
    \item Recomendar arranjos institucionais específicos: a passagem de ``correlação observada entre formação econômica e percepção em $n=184$'' para ``portanto são necessários órgãos técnicos independentes'' requer, no mínimo, (a) replicação em amostra probabilística; (b) identificação causal do efeito da formação econômica; (c) evidência de que as diferenças observadas refletem acurácia superior e não diferenças de valores ou preferências.
\end{enumerate}

\subsection*{Agenda de pesquisa}

Para que os padrões aqui documentados se traduzam em conhecimento acumulável, seria necessário: (i) estudo confirmatório com $n_{\texttt{econ}} \geq 120$ e amostragem probabilística ou estratificada por região; (ii) experimento informacional à la \citeonline{kuziemko2015egalitarian}, testando o efeito de apresentar informações econômicas específicas sobre crenças mensuradas antes e depois; (iii) análise de heterogeneidade regional, dado que o Brasil apresenta disparidades de desenvolvimento e composição ideológica que podem moderar os padrões observados.

\citeonline{hayek_knowledge_use} identificou o problema do conhecimento disperso: nenhum planejador central pode agregar o conhecimento local fragmentado. A dominância ideológica documentada acrescenta uma camada ao problema — o conhecimento é também filtrado por predisposições identitárias. Esse diagnóstico, que \textit{os dados suportam com as qualificações acima}, é compatível com a literatura de public choice \cite{buchanan1962calculus} e de governança de bens comuns \cite{ostrom1990governing}. Mas a passagem da descrição associacional para a prescrição institucional exige um degrau de identificação causal que este estudo exploratório não alcança.

\section{Limitações e Agenda Futura}
\label{sec:limitacoes-v2}

A limitação principal é $n_{\texttt{econ}} = 17$. Um estudo confirmatório que busque detectar OR $= 2{,}0$ com 80\% de poder e $\alpha = 0{,}05$ requereria $n_{\texttt{econ}} \geq 120$, o que implica amostra total de $\approx 1.300$ respondentes com a mesma proporção de 9\%. Este estudo é, portanto, \textit{exploratório com rigor metodológico} — não confirmatório.

A amostragem não probabilística limita a generalização: os resultados descrevem associações condicionais dentro do quadro amostral, não prevalências populacionais. Os modelos estimam associações, não efeitos causais. A formação econômica avançada é endógena: quem escolhe fazer mestrado ou doutorado em Economia difere sistematicamente de quem não escolhe, em dimensões que nenhum conjunto de covariáveis observadas consegue capturar inteiramente.

\paragraph{Sobre a composição ideológica do grupo econ=1 e a endogeneidade do déficit.}
Um ponto metodológico específico merece atenção explícita. Dos 14 economistas com espectro político válido, 57\% se declaram de esquerda e apenas 29\% de direita (Tabela~\ref{tab:crosstab_econ_esp}). No item \textit{déficit federal grande}, economistas concordam menos com a afirmação do que o grupo de controle ($\Delta\bar{Y} = -0{,}46$). Uma hipótese alternativa plausível é que esse padrão reflita a composição ideológica do grupo econ=1 (respondentes de esquerda tendem a ser mais céticos quanto à gravidade do déficit) e não a formação técnica em si. O modelo controla por \texttt{espectro} individualmente, o que atenua (mas não elimina) essa confusão: a correlação entre formação e ideologia não observada residual persiste. Em outras palavras, não é possível descartar a hipótese alternativa com este design — é exatamente por isso que o estimando $\hat{\beta}_{\texttt{econ}}$ não pode ser interpretado causalmente. Um estudo com instrumento exógeno para a formação econômica (e.g., variação de oferta de vagas em programas de pós-graduação por região) seria necessário para separar o efeito da formação do efeito de seleção ideológica.

A agenda futura natural inclui: (i) estudo confirmatório com $n_{\texttt{econ}} \geq 120$, preferencialmente com amostragem probabilística ou estratificada por região; (ii) experimento informacional à la \citeonline{kuziemko2015egalitarian}, testando o efeito causal de apresentar informações econômicas específicas sobre crenças; (iii) análise de heterogeneidade regional, dado que o Brasil apresenta disparidades de desenvolvimento e de composição ideológica que podem moderar os efeitos observados; (iv) estudo longitudinal que acompanhe a evolução das percepções com e sem intervenção educacional.

\section{Síntese Final}
\label{sec:sintese-final}

Os resultados delineiam o seguinte quadro: a posição ideológica autodeclarada está associada à maior parte da variação observada nas percepções econômicas desta amostra; a formação econômica avançada está associada a diferenças condicionais em um subconjunto específico de itens — precisamente aqueles com maior componente factual. Esses são padrões associativos, não efeitos causais identificados.

Esse diagnóstico não é pessimista; é realista. E o realismo tem consequências práticas. Primeiro, a comunicação baseada em evidência precisa reconhecer que a batalha não se dá apenas no terreno dos fatos — ela se dá também no terreno das molduras narrativas. Segundo, a educação econômica tem valor, mas seu valor é condicional ao design: currículos que cultivam humildade intelectual, lógica de identificação e hábito de \textit{disconfirmation} produzem economistas mais robustos à autocognição cultural do que currículos que apenas transmitem modelos \cite{franco2022cartas}. Terceiro, instituições que amortizam erros previsíveis são o complemento necessário da educação, não seu substituto.

Newman entendia que a educação liberal não é transmissão de informação, mas formação do intelecto: a capacidade de ver conexões, suspeitar das próprias intuições, reconhecer os limites do próprio juízo \cite{newman2020ideia}. No campo econômico, isso significa aprender a procurar custos invisíveis, efeitos de segunda ordem e as consequências não intencionais das políticas que parecem óbvias. Não é modelagem matemática; é disciplina do olhar.

À luz do princípio popperiano, as conclusões são provisórias e abertas à refutação \cite{popperlogic}. Elas apontam um caminho: regras que desincentivem erros previsíveis; subsidiariedade que devolva decisões ao nível onde o conhecimento é mais rico; e educação que reduza a demanda por más políticas formando juízo, não apenas transmitindo modelos. Quando o erro é informacional, educa-se; quando é identitário, desenham-se instituições que limitem seu dano. Distinguir um do outro — item a item, contexto a contexto — é a tarefa que a econometria comportamental pode ajudar a cumprir.
